% QUESTAO
Dadas duas matrizes $A$ e $B$ de mesma ordem, denominamos diferença entre $A$ e $B$, indicada por $A-B$, a matriz $C$ obtida ao calcularmos a adição de $A$ com o oposto de $B$, ou seja, $A-B=A+(-B)=C$. Assim, sejam $A = \left[\begin{array}{cc}
-1 & 3 \\
2 & -8
\end{array}\right]$ e $B = \left[\begin{array}{cc}
2 & -1 \\
3 & 0
\end{array}\right]$. Calcule a diferença $B-A$.

% ALTERNATIVAS
$B-A = \left[\begin{array}{cc} 3 & -4 \\ 1 & 8 \end{array}\right]$
$B-A = \left[\begin{array}{cc} -3 & 4 \\ -1 & -8 \end{array}\right]$
$B-A = \left[\begin{array}{cc} 6 & -8 \\ 2 & -16 \end{array}\right]$
$B-A = \left[\begin{array}{cc} \frac{3}{2} & -2 \\ \frac{1}{2} & -4 \end{array}\right]$
$B-A = \left[\begin{array}{cc}  1 & -8 \\  3 & -4 \end{array}\right]$


